\documentclass[11pt]{article}
\usepackage[a4paper,margin=28mm]{geometry}
\usepackage{amsmath,amssymb,amsthm}
\usepackage[hidelinks]{hyperref}
\newtheorem{proposition}{Proposition}
\newtheorem{lemma}[proposition]{Lemma}
\newcommand{\gen}[1]{\langle #1\rangle}
\title{Panels in the coset construction for a $C^+$-group}
\author{Human Compression project}
\date{}
\begin{document}
\maketitle
\begin{abstract}
We distinguish a rank restriction missing from the literal statement of
Corollary 4.9 of Leemans--Tranchida from the unresolved assertion in higher
rank. We also give an exact criterion for geometry and firmness in rank
four, using products of the four maximal parabolics. This reduces a
possible finite counterexample to a small certificate without an
automorphism computation. The criterion is checked on four previously
acquired markings; none is a counterexample in rank four.
\end{abstract}

\section{The rank restriction}
We use the final published edition of \cite{LT}, printed pages 493--494
and 497. Put $\alpha_0=1$ and
\[
 P_J=\gen{\alpha_i^{-1}\alpha_j:i,j\in J}.
\]
The $C^+$ intersection condition is
$P_J\cap P_K=P_{J\cap K}$ for $|J|,|K|\geq2$.
Construction 3.2 takes $H_i=P_{I\setminus\{i\}}$ and forms the incidence
system of right cosets $H_i g$, with incidence defined by nonempty
intersection. A geometry means that every flag extends to a flag
containing one object of each type. Connectedness is a separate condition.

\begin{proposition}
For every nontrivial cyclic group $G=\gen{a}$, the rank-two marking
$R=\{a\}$ satisfies the stated $C^+$ condition and is independent.
Its associated coset system is a geometry, but is not firm.
\end{proposition}
\begin{proof}
Here $I=\{0,1\}$, $\alpha_0=1$, and $\alpha_1=a$.
The only subset of $I$ with at least two elements is $I$ itself, so the
required intersection equality is $G\cap G=G$. Independence says
$a\notin\gen{\varnothing}=1$, which holds.
Both maximal parabolics are trivial, since their defining index sets
are singletons. Each type therefore has one object for each $g\in G$.
The type-zero object labelled $g$ is incident to the type-one object
labelled $h$ exactly when $g=h$. The geometry consists of disjoint edges.
Every flag extends to an edge, but every one-object panel has exactly
one completion. This is the failure of firmness.
\end{proof}

No lower rank bound is stated in the displayed $C^+$ definition,
Construction 3.2 or Corollary 4.9. The proof of Lemma 4.8, however,
chooses an index $k$ distinct from two prescribed indices $i,j$.
Consequently this example identifies a missing rank restriction in the
literal wording of both Lemma 4.8 and Corollary 4.9: replacing one face
of a matching edge gives a nonincident pair, not an adjacent chamber.
Section 2.1, printed page 492, imposes neither connectedness nor a lower
rank bound on a geometry, and its Proposition 2.1 permits every positive
number of types. The later definitions do not add a global restriction.
The counterexample does not decide the intended statement for rank at
least three. In rank three the desired conclusion follows from pair
transitivity and the canonical adjacent chambers, as in the existing
project proof of Proposition 4.11. For rank at least four, the remaining
problem is to supply a neighbour for a panel outside the orbit of a
canonical panel. The accepted counterexample to the principal
residual-connectedness theorem is thin, and does not solve this problem.

\section{A criterion in rank four}
The next reduction applies to arbitrary subgroups of a finite group;
the $C^+$ condition is needed only for the last description of the
stabilizer. In particular, no transitivity on triple flags is assumed.

\begin{lemma}
Let $G$ be finite, let $I=\{0,1,2,3\}$, and let $H_i\leq G$.
For each missing type $d$, choose an ordering $(a,b,c)$ of
$I\setminus\{d\}$ and put
\[
 X_c=H_cH_a\cap H_cH_b,\qquad
 Y_d(x)=H_dH_a\cap H_dH_b\cap H_dH_cx \quad(x\in X_c).
 \tag{1}
\]
The coset system is a geometry if and only if $Y_d(x)\ne\varnothing$
for all four choices of $d$ and all $x\in X_c$. It is a firm geometry
if and only if
\[
 |Y_d(x)|\geq 2|H_d| \quad\text{for all these }d,x.
 \tag{2}
\]
It suffices in either test to take one representative for each orbit of
$K=H_a\cap H_b$ on the right cosets $H_cx$ contained in $X_c$.
For a rank-four $C^+$ marking,
$K=P_{\{c,d\}}=\gen{\alpha_c^{-1}\alpha_d}$ is cyclic.
\end{lemma}
\begin{proof}
Every incident pair can be taken to $H_a,H_b$ by right multiplication:
choose an element $z$ in the intersection of its two cosets and
multiply both cosets by $z^{-1}$. Every pair, every object and the empty
flag therefore extend to a canonical chamber after the same
normalization. The only possible obstruction to being a geometry is a
triple flag with no completion at the fourth type.

After normalizing the types $a,b$ of such a triple, its last object is
$H_cx$. It intersects both $H_a$ and $H_b$ precisely when $x\in X_c$.
A type-$d$ coset $H_dy$ completes the triple precisely when $y$ belongs
to each of the three sets defining $Y_d(x)$. Each such set, and hence
their intersection, is a union of right cosets of $H_d$. The number of
completions is therefore exactly $|Y_d(x)|/|H_d|$. This proves the
geometry criterion and (2). Once all panels have at least two
completions, every flag of corank one has a residue with at least two
objects, which is precisely firmness.

Changing $x$ to $h x$ with $h\in H_c$ leaves $Y_d(x)$ unchanged.
Right multiplication by $k\in K$ fixes $H_a,H_b$ and sends the panel
with third object $H_cx$ to the panel with third object $H_cxk$.
It bijects their completions; equivalently $Y_d(xk)=Y_d(x)k$.
Thus one representative from each stated orbit suffices.
Finally the intersection condition gives
$H_a\cap H_b=P_{I\setminus\{a,b\}}$.
The subgroup for a two-element index set is generated by either one of
the mutually inverse differences, giving the cyclic description.
\end{proof}

\section{Finite checks and remaining scope}
The accompanying program checks the rank-two example for $G=C_2$ by
enumerating all flags: there are four objects, two chambers, and four
panels, each with one completion. This check is separate from the
uniform written proof for all nontrivial cyclic groups.

For rank four the program reuses the four stored $\operatorname{PSL}(2,7)$
markings from the earlier diagnostic. It checks the multiplication
table against determinant-one matrices modulo their central signs,
regenerates all sixteen index subgroups, and verifies their 256
intersection equalities. It checks (1) against literal coset
intersections for every normalized panel, then partitions the panels
into the stated stabilizer orbits. No marking is searched for anew.
Each of the four geometries has five completions at every panel.
The existing projective-line description gives the same conclusion:
four typed copies of an eight-point set are incident by inequality.

For a future fixed finite group, a proposed counterexample must pass
the nondegeneracy and full intersection checks, have no zero entry in
(1), and have at least one entry equal to $|H_d|$. A singleton panel
without the first of these two completion requirements would not
refute Corollary 4.9: the coset system might fail to be a geometry.
This packet does not produce a counterexample or a general proof in
rank at least four, and it gives no whole-original, conjecture,
formalization or compression credit.

The $C^+$ construction and canonical neighbour formula are due to the
authors of \cite{LT} and their cited predecessors. The previously
written chamber-component argument and the four stored markings are
reused. The present additions are the explicit rank boundary and the
panel-counting reduction; no claim of historical novelty for the
elementary coset calculation is made.

\begin{thebibliography}{1}
\bibitem{LT} D.~Leemans and P.~Tranchida,
\emph{On residual connectedness in chiral geometries},
Algebraic Combinatorics \textbf{4} (2021), 491--499,
\href{https://doi.org/10.5802/alco.162}{doi:10.5802/alco.162}.
\end{thebibliography}
\end{document}
